\section{Background and Motivation}
\label{sec:background}

\subsection{Differential Privacy and Its Accounting}
\label{sec:bg-dp}

A randomized mechanism $M$ satisfies $(\eps,\delta)$-differential
privacy~\cite{dwork2006differential} if for
every pair of neighbouring datasets $D,D'$ and every event $S$,
$\Pr[M(D)\in S]\le e^{\eps}\Pr[M(D')\in S]+\delta$. The definition is
per-release, but its power comes from composition: the guarantee for a transcript
follows from the guarantees of its parts. That is what lets a platform promise a
subject a bound on \emph{total} leakage rather than on each answer, and it is why
the promise needs bookkeeping. Whatever its interpretation---one of the variants surveyed
in~\cite{desfontaines2019differential}---the quantity to bound is cumulative.

Composition bounds differ in how tightly they track it. R\'enyi
DP~\cite{mironov2017renyi} gives sharp constants for a fixed number of adaptively
chosen mechanisms, at the cost of a discrete order grid: the accountant evaluates
the grid and converts the result to $(\eps,\delta)$.
Zero-concentrated DP~\cite{bun2016concentrated} bounds the R\'enyi divergence of
order $1+\alpha$ by $\rho\alpha$ for all $\alpha>1$, composes by adding $\rho$,
and converts in closed form,
\begin{equation}
\eps = \rho + 2\sqrt{\rho\ln(1/\delta)} ,
\label{eq:zcdp}
\end{equation}
with inverse $\rho = \bigl(\sqrt{\ln(1/\delta)+\eps}-\sqrt{\ln(1/\delta)}\bigr)^{2}$.
A Gaussian mechanism with $\ell_2$-sensitivity $\Delta$ and noise scale $\sigma$
is $\rho$-zCDP with $\rho=\Delta^{2}/(2\sigma^{2})$, so
$\sigma=\Delta/\sqrt{2\rho}$ for a target $\rho$.
Equation~\eqref{eq:zcdp} is not the tightest known conversion---$f$-DP and
Gaussian DP sharpen composition for Gaussian mechanisms~\cite{dong2021rejoinder}---but
it is monotone, invertible and cheap, which is what an admission decision needs.
Section~\ref{sec:eval-backend} shows that an R\'enyi accountant cannot calibrate a
target of $\eps{=}0.02$ at all on a practical order grid.

\subsection{The Market Setting}
\label{sec:bg-market}

We consider a market in which subjects contribute records to one or more
platforms, and tenants buy answers over those records. A permissioned ledger
records contributions, query charges and settlements, so that a tenant can verify
that an answer corresponds to a committed query and that a charge was not
silently dropped. Trustless trading designs already provide this for
delivery~\cite{liu2025yotta}, and provenance mechanisms extend it to derived
artifacts~\cite{zhu2025auditing}. What the ledger does not provide is a statement
about a subject's remaining budget, because that depends on queries answered
elsewhere and on charges still in flight.

Three properties of the setting shape the design. The subject population is large
and heavy-tailed: most subjects appear in a handful of queries, a small minority
in thousands, and it is that minority whose budgets bind. Tenants are numerous
and mutually distrustful, and a tenant's incentive is to maximise the information
it extracts per unit of budget, including under several identities. Queries
arrive continuously, so the enforcement decision sits on the critical path of
every answer the platform sells and has a latency budget of microseconds to
milliseconds.

\subsection{Approximate Counting for Thresholds}
\label{sec:bg-sketch}

A Count-Min sketch stores a $d\times w$ matrix of counters and updates, for key
$x$ and weight $v$, one cell per row at the positions given by $d$ independent
hash functions~\cite{fusy2023countmin}. The estimate is the minimum of the $d$
cells, so collisions can only inflate it. Conservative update strengthens the
guarantee: instead of adding $v$ to every cell it computes the current minimum
$m$ and raises each cell to at most $m+v$, which never lowers an estimate and
reduces overestimation in
practice~\cite{fusy2022phase,fusy2023countmin}. A space-saving summary over the
same stream retains the $k$ largest keys exactly and evicts the smallest, which
is the standard way to recover heavy hitters. Recent work replaces the hash
functions with a learned model~\cite{nishishita2025optimized} and revisits the
row-count trade-off~\cite{fusy2023countmin}; both optimise estimation error,
whereas our requirement is weaker and, in one respect, stricter: the estimate
must never fall below the truth, because it gates a safety property rather than a
reported statistic.

\subsection{Service-Level Objectives}
\label{sec:bg-slo}

An SLO is a predicate over a measurement window, not over a lifetime total. The
distinction matters because a subject's tolerance is typically expressed as ``no
more than $\eps$ in any thirty days'', which is a sliding window and therefore a
different object from the cumulative $\eps$ that composition bounds. Enforcing only
the total exhausts a budget in the first hour; enforcing only the window lets the
total drift without bound. \sys\ maintains both: a cumulative sketch that never forgets, and a ring of epoch
sketches for windowed queries, with the SLO on the window and a separate hard cap on
the total.
